\documentclass[11pt]{ctexart}
\usepackage[margin=2.5cm]{geometry}
\usepackage{amsmath,amssymb}
\usepackage{booktabs,tabularx}
\usepackage[authoryear,round]{natbib}
\usepackage[colorlinks=true,linkcolor=blue,citecolor=blue,urlcolor=blue]{hyperref}
\usepackage{array}


\newcommand{\E}{\mathbb{E}}
\newcommand{\F}{\mathcal{F}}

\begin{document}

\setcounter{section}{1}
\section{基础凸随机}

\subsection{设定}

$f$ 凸、$L$-光滑，$x^\star\in\arg\min f$，$S:=\arg\min f$，$D:=\|x_0-x^\star\|$，$\Delta_t:=f(x_t)-f^\star$。SGD \citep{robbins1951} 为
\[
x_{t+1}=x_t-s_t\bigl(\nabla f(x_t)+e_t\bigr),
\qquad
\E_t[e_t]=0,
\quad
\E_t\|e_t\|^2\le\sigma^2,
\]
其中 $\E_t[\cdot]:=\E[\cdot\mid\F_t]$。

\subsection{一步不等式}

当 $s_t\le 1/(2L)$ 时，
\begin{equation}
\boxed{\;
\E_t\|x_{t+1}-x^\star\|^2+s_t\bigl(f(x_t)-f^\star\bigr)
\le \|x_t-x^\star\|^2+s_t^2\sigma^2 .
\;}
\tag{SGD}
\end{equation}
三个要素：条件无偏性消掉交叉项；凸性给出函数值间隙；光滑性给出 $\|\nabla f(x)\|^2\le 2L\bigl(f(x)-f^\star\bigr)$ \citep[Theorem~2.1.5]{nesterov2004}。
% 推导见附录~\ref{app:sgd}，
框架同 \citet{nemirovski2009} 与 \citet{bottou2018}。

\subsection{偏差--方差界}

对 (SGD) 从 $0$ 到 $k-1$ 求和，再用 Jensen 不等式。记
\[
S_k:=\sum_{t<k}s_t,
\qquad
\widetilde x_k:=\frac1{S_k}\sum_{t<k}s_tx_t ,
\]
得到 \citep{nemirovski2009,bottou2018}
\begin{equation}
\boxed{\;
\E[f(\widetilde x_k)-f^\star]
\le
\underbrace{\frac{D^2}{S_k}}_{\text{偏差}}
+\underbrace{\frac{\sigma^2\sum_{t<k}s_t^2}{S_k}}_{\text{方差}} .
\;}
\tag{9}
\end{equation}
这是本部分的核心公式。第 1 部分中 JOTA 的界与它同构，只是把 $S_k$ 换成 $s_kt_k^2$。

\subsection{由 (9) 得到的速率}

\begin{center}
\begin{tabularx}{\linewidth}{@{}>{\raggedright\arraybackslash}X l >{\raggedright\arraybackslash}X@{}}
\toprule
步长 & 期望速率（加权平均点） & 说明 \\
\midrule
常数 $s\asymp\min\{1/L,\,D/(\sigma\sqrt T)\}$，已知 $T$ & $O\bigl(LD^2/T+\sigma D/\sqrt T\bigr)$ & 平衡 $D^2/(sT)$ 与 $s\sigma^2$ \citep{nemirovski2009} \\
$s_t\asymp t^{-p}$，$p<1/2$ & $O(k^{-p})$ & 方差主导 \\
$p=1/2$ & $O(\log k/\sqrt k)$ & 与 $T$ 无关的最好选择 \\
$1/2<p<1$ & $O(k^{-(1-p)})$ & $\sum_t s_t^2<\infty$ \citep{robbins1951} \\
\bottomrule
\end{tabularx}
\end{center}
后三行是把 $s_t\asymp t^{-p}$ 直接代入 (9) 的计算。

\subsection{下界：加速只能改善偏差项}

在凸光滑、方差有界的随机一阶 oracle 模型下，minimax 下界为
\[
\Omega\Bigl(\frac{LD^2}{T^2}+\frac{\sigma D}{\sqrt T}\Bigr).
\]
\begin{itemize}
  \item 偏差项的下界来自确定性一阶方法的下界 \citep[Theorem~2.1.7]{nesterov2004}；
  \item 噪声项的下界来自信息论论证 \citep{nemirovski1983,agarwal2012}；
  \item 加速方法 AC-SA 同时达到两项 \citep{lan2012}。
\end{itemize}
所以 SGD 在噪声项上已经最优，加速只把偏差项从 $LD^2/T$ 改进到 $LD^2/T^2$。当 $k\to\infty$ 时，渐近阶由 $\sigma D/\sqrt k$ 决定：\textbf{持续噪声下不存在渐近意义的加速。}

\subsection{几乎必然分析的入口}

若 $\sum_t s_t=\infty$、$\sum_t s_t^2<\infty$ \citep{robbins1951}，对 (SGD) 应用 Robbins--Siegmund 引理 \citep{robbins1971}，a.s.\ 有：
\begin{enumerate}
  \item 对每个 $x^\star\in S$，$\|x_t-x^\star\|$ 收敛（因此轨道有界）；
  \item $\sum_t s_t\bigl(f(x_t)-f^\star\bigr)<\infty$。
\end{enumerate}
由第 2 条和 Jensen 不等式，加权平均点满足
\[
f(\widetilde x_k)-f^\star=O\bigl(1/S_k\bigr)\quad\text{a.s.}
\]
取 $s_t\asymp t^{-p}$，$p\in(1/2,1)$，即为 $O(k^{-(1-p)})$，\textbf{任意接近但达不到} $k^{-1/2}$。换用适当的平均方式可以把 $O$ 改进为 $o$ \citep{sebbouh2021}。

由第 1 条加上聚点论证，迭代点 a.s.\ 收敛到某个随机极小点 \citep{opial1967,combettes2015}。概率一事件依赖于 $x^\star\in S$，而 $S$ 通常不可数，需要取可数稠密子集处理这一量词问题 \citep{combettes2015}。末次迭代 $x_t$ 的 a.s.\ 速率需要另外的论证 \citep{sebbouh2021,liu2022}。

\textbf{与第 1 部分的对应：}Robbins--Siegmund 对应``能量收敛''，$\sum_t s_t\Delta_t<\infty$ 对应 (K1$^+$) 的余量，Opial 对应迭代点收敛。

% \appendix
% \section{一步不等式的推导}\label{app:sgd}

% 由 $x_{t+1}=x_t-s_t\bigl(\nabla f(x_t)+e_t\bigr)$ 展开距离平方：
% \begin{equation}
% \begin{aligned}
% \|x_{t+1}-x^\star\|^2
% ={}&\|x_t-x^\star\|^2
% -2s_t\langle\nabla f(x_t),x_t-x^\star\rangle
% -2s_t\langle e_t,x_t-x^\star\rangle\\
% &+s_t^2\|\nabla f(x_t)+e_t\|^2 .
% \end{aligned}
% \tag{A.1}
% \end{equation}
% 取条件期望。由 $\E_t[e_t]=0$，线性交叉项消失：
% \begin{equation}
% \E_t\|x_{t+1}-x^\star\|^2
% =\|x_t-x^\star\|^2
% -2s_t\langle\nabla f(x_t),x_t-x^\star\rangle
% +s_t^2\E_t\|\nabla f(x_t)+e_t\|^2 .
% \tag{A.2}
% \end{equation}
% 由凸性，
% \begin{equation}
% \langle\nabla f(x_t),x_t-x^\star\rangle\ge f(x_t)-f^\star=\Delta_t .
% \tag{A.3}
% \end{equation}
% 由凸性和 $L$-光滑性 \citep[Theorem~2.1.5]{nesterov2004}，
% \begin{equation}
% \|\nabla f(x_t)\|^2\le 2L\Delta_t .
% \tag{A.4}
% \end{equation}
% 条件无偏性还给出
% \begin{equation}
% \E_t\|\nabla f(x_t)+e_t\|^2
% =\|\nabla f(x_t)\|^2+\E_t\|e_t\|^2
% \le 2L\Delta_t+\sigma^2 .
% \tag{A.5}
% \end{equation}
% 将 (A.3)--(A.5) 代入 (A.2)：
% \begin{equation}
% \E_t\|x_{t+1}-x^\star\|^2
% \le\|x_t-x^\star\|^2-s_t(2-2Ls_t)\Delta_t+s_t^2\sigma^2 .
% \tag{A.6}
% \end{equation}
% 当 $s_t\le1/(2L)$ 时 $2-2Ls_t\ge1$，即得 (SGD)：
% \begin{equation}
% \E_t\|x_{t+1}-x^\star\|^2+s_t\Delta_t\le\|x_t-x^\star\|^2+s_t^2\sigma^2 .
% \tag{A.7}
% \end{equation}

\bibliographystyle{plainnat}
\bibliography{refs}

\end{document}
